%!TEX root = ../main.tex
% =============================================================================
%  三、数据预处理与特征分析      建议 2–3 页
%  评审在问：数据你读懂了吗？
%  本节必须出现：① 时标与对齐结论  ② 附件 1 身份核查（= 全年逐时段均值，非预报）
%                ③ 相关性的尺度拆解  ④ 预报误差分档（全样本 / 白天双口径）
%  禁止：相关性写成因果（"导致""引起""造成"一律换成"共动""同向变动"）
% =============================================================================

\section{数据预处理与特征分析}

\subsection{数据来源与基本规模}

本文使用的四个附件及其用途如表~\ref{tab:data} 所示。所有时间序列均按 \textbf{10~min} 粒度、
以\textbf{时段末端}为标签对齐，全天共 144 个时段。

\begin{table}[htbp]
  \centering
  \caption{附件数据规模、结构与本文用途（数据行数，不含表头）}
  \label{tab:data}
  \footnotesize
  \begin{tabularx}{\linewidth}{C{1.5cm} L{3.4cm} C{2.4cm} L{3.0cm} Y}
    \toprule
    附件 & 内容 & 数据行数 & 时间范围 & 本文用途 \\
    \midrule
    附件 1 & 电价、小区负载、光伏发电"预测"功率 & 144 & 一天（典型日） & 问题 1 的输入；作为"全年平均日"使用 \\
    附件 2 & 小区负载、光伏发电实际功率 & \(365\times144\) & 2025-01-01 至 12-31 & 问题 2--4 的实际值，用于结算与回测 \\
    附件 3 & 光伏发电功率预报（整点） & \(365\times4\times24\) & 同上，每日 4 次发布 & 问题 3、4 的日前与滚动预报输入 \\
    附件 4 & 电价 & \(365\times144\) & 同上 & 问题 4 的波动电价输入 \\
    \bottomrule
  \end{tabularx}
\end{table}

\subsection{时标语义与对齐}

\begin{guideblock}
这是\textbf{结果整体错位的头号来源}，必须在正文写清三件事：
① 标签的语义是"时段末端"；② 附件 3 的"预报 \(k\) 小时"落在哪个小时；
③ 每天恰有 144 个时段，\textbf{不得丢弃最后一列} \(0{:}00{+}1\)。
\end{guideblock}

附件 1、2、4 的时段标签为 \(0{:}10,0{:}20,\dots,23{:}50,0{:}00{+}1\)（后者表示次日
\(0{:}00\)）。本文约定：

\begin{itemize}
  \item 标签表示\textbf{时段末端}：第 \(i\) 个数据点代表区间
        \([(i-1)\Delta t,\ i\Delta t)\)，\(i=1,\dots,144\)，\(\Delta t=1/6\)~h；
  \item 每一行的 144 个值\textbf{无缝铺满} \(0{:}00\) 至 \(24{:}00\)，
        最后一列 \(0{:}00{+}1\) 是当日 \(23{:}50\)--\(24{:}00\) 的量，\textbf{不得丢弃}——
        否则每处合计偏低 \(1/144\approx 0.6944\%\)。
\end{itemize}

附件 3 的列标签为"预报 1 小时"至"预报 24 小时"，本文约定
\textbf{预报 \(k\) 小时对应发布时刻之后的第 \(k\) 个小时}，即 \((h_0+k-1,\ h_0+k]\)
（\(h_0\in\{0,6,12,18\}\)）。以 2025-01-01 \(0{:}00\) 发布的"预报 7 小时"为例，
其值对应 \(6{:}00\)--\(7{:}00\) 的光伏预报功率；若误按 \((h_0+k,\ h_0+k+1]\) 对齐，
则\textbf{全部预报整体错 1 小时}。\(18{:}00\) 的预报中 \(k=7,\dots,24\) 落在
\textbf{次日} \(0{:}00\)--\(17{:}00\)，本文用绝对时间轴映射
（目标小时序号 \(=d\times24+h_0+k-1\)）处理跨日配对，越界配对丢弃并记录数量 \ph~条。
附件 3 的日期列为稀疏填充，需前向填充（\texttt{ffill}）后断言每天恰有 4 行。
表~\ref{tab:timeruler} 给出题面时段与附件标签的对照。

\begin{table}[htbp]
  \centering
  \caption{题面表 1 的时段写法与附件时段标签对照（示例）}
  \label{tab:timeruler}
  \small
  \begin{tabular}{cccc}
    \toprule
    题面写法 & 附件标签 & 对应区间 & 时段序号 \(i\) \\
    \midrule
    \(10{:}00\)--\(10{:}10\) & \(10{:}10\) & \([10{:}00,10{:}10)\) & 61 \\
    \(12{:}00\)--\(12{:}10\) & \(12{:}10\) & \([12{:}00,12{:}10)\) & 73 \\
    \(14{:}00\)--\(14{:}10\) & \(14{:}10\) & \([14{:}00,14{:}10)\) & 85 \\
    \(16{:}00\)--\(16{:}10\) & \(16{:}10\) & \([16{:}00,16{:}10)\) & 97 \\
    \(20{:}00\)--\(20{:}10\) & \(20{:}10\) & \([20{:}00,20{:}10)\) & 121 \\
    \(23{:}50\)--\(24{:}00\) & \(0{:}00{+}1\) & \([23{:}50,24{:}00)\) & 144 \\
    \bottomrule
  \end{tabular}
\end{table}

\begin{guideblock}
对齐自检（必须写进代码并在本节报告）：以 2025-01-01 的 \(0{:}00\) 预报"预报 7 小时"
为例，核对其值应对应 \(6{:}00\)--\(7{:}00\) 的光伏预报；再校验
"全天汇总 = \(\sum\) 144 个时段值"（相对误差 \(\le 10^{-4}\)）——
这一条断言能同时挡住"丢最后一列"（偏低 0.6944\%）与"错位一个槽"两类错误。
\end{guideblock}

\subsection{附件 1 的身份核查}

附件 1 第 4 列的表头写作"光伏发电\textbf{预测}功率"，但经核查，其数值等于附件 2 全年
逐时段实际光伏功率的\textbf{逐时段均值}（相关系数 \ph，逐时段偏差最大值 \ph~kW）；
其负载与电价同样按全年逐时段均值理解。因此附件 1 代表\textbf{全年平均日（典型日）}，
\textbf{不是}日前预报，\textbf{不可用作问题 2--4 的验证集}（否则构成循环论证）。
问题 1 因此是一个\textbf{确定性优化问题}，其最优值是问题 2--4 的对照基准。

\begin{guideblock}
本小节是"数据你读懂了吗"最直接的证据，\textbf{不要省掉}。
核查结论必须由本队代码算出并给出相关系数与最大偏差，不能只写定性描述。
\end{guideblock}

\subsection{电价与负载的结构特征}

\begin{guideblock}
本节保留三样真正会影响模型的东西：
① 峰谷价差是否超过盈亏平衡阈值 \(1/(\eta_c\eta_d)\approx 1.235\)（决定"能不能套利"）；
② 光伏盈余时段的规模（决定"容量是不是瓶颈"）；
③ 相关性的\textbf{尺度依赖}（决定"要不要建模联合分布"）。
凡是与模型无关的描述性统计，一律删掉。
\end{guideblock}

按附件 4 统计，电价均值 \ph~元/kWh、极值区间 \ph~至 \ph~元/kWh，呈
"夜间低谷—上午高峰—午间回落—傍晚尖峰"形态；按附件 2 统计，负载均值 \ph~kW、
光伏均值 \ph~kW，净负荷（\(L_t-PV_t\)）为负即光伏盈余，其时段占比 \ph。

套利是否有利取决于峰谷价差能否覆盖往返损耗：在电价 \(p_1\) 处充电、\(p_2\) 处放电时，
单位电量净收益为 \(p_2\eta_c\eta_d-p_1\)，故\textbf{盈亏平衡价比阈值为
\(1/(\eta_c\eta_d)\approx 1.235\)}。本文实测谷峰价比为 \ph（时段划分：\ph，须写明），
高于阈值；而盈余功率超过充电功率上限 \(\overline P=5000\)~kW 的时段数为 \ph~个，
说明\textbf{价差充裕但容量受限}——这正是本文核心矛盾所在：
不在于"能不能套利"，而在于"如何把有限的储能容量分配到最优时段"。

电价与负载的相关系数具有\textbf{尺度依赖性}，表~\ref{tab:corr} 按三个尺度分别报告。

\begin{table}[htbp]
  \centering
  \caption{电价与负载、净负荷、光伏的相关系数（三个尺度，样本量 \(n\) 已注明）}
  \label{tab:corr}
  \small
  \begin{tabular}{lccc}
    \toprule
    变量对 & 池化（全部时段） & 日级（逐日先算再合并） & 剔除日内形状后 \\
    \midrule
    电价 --- 负载 & \ph & \ph & \ph \\
    电价 --- 净负荷 & \ph & \ph & \ph \\
    电价 --- 光伏 & \ph & \ph & \ph \\
    \bottomrule
  \end{tabular}
\end{table}

\begin{guideblock}
三条纪律：① 必须写明是\textbf{哪一对变量}——"电价—净负荷"与"电价—负载"的数值不同，
不写清楚会让论文自相矛盾；② 三个尺度的结论可以不同，这本身就是发现，要如实报告；
③ \textbf{只能写"共动""同向变动""结构上重合"，不得写"导致""引起""造成"}——
价格与负载由同一作息节律驱动，属共动而非因果。
\end{guideblock}

\subsection{光伏预报误差的统计刻画}

\begin{guideblock}
必须\textbf{双口径并报}（全样本 / 白天且实际功率大于阈值），并写明筛选条件与样本量 \(n\)。
只报一个 MAE 会被质疑样本构成假象。提前期效应必须用\textbf{配对}设计报告，
不得只写"N 倍"而不给配对方式与 \(n\)。
\end{guideblock}

本文将附件 3 的全部预报值与附件 2 的对应实际值配对，按\textbf{提前期}（1--24 小时）
与\textbf{目标整点}两个维度统计误差。考虑到夜间光伏功率接近 0、相对误差无意义，
本文同时报告\textbf{全样本口径}与\textbf{白天口径}（实际功率 \(>100\)~kW），
结果见表~\ref{tab:fcsterr}。

\begin{table}[htbp]
  \centering
  \caption{光伏预报误差统计（两个口径；MAE、RMSE、bias 单位均为 kW，相对误差为 MAE 与白天实际均值的比值）}
  \label{tab:fcsterr}
  \small
  \begin{tabular}{lcccc}
    \toprule
    口径 & 样本量 \(n\) & MAE & RMSE & bias \\
    \midrule
    全样本 & \ph & \ph & \ph & \ph \\
    白天（实际 \(>100\)~kW） & \ph & \ph & \ph & \ph \\
    \bottomrule
  \end{tabular}
\end{table}

误差随提前期的变化采用\textbf{同一目标整点配对}的方式比较：
以目标整点为配对键，比较"提前 \(k_1\) 小时"与"提前 \(k_2\) 小时"的预报误差，
配对样本量 \(n=\ph\)，误差比值为 \ph（配对方式：\ph）。
本文以\textbf{白天口径}作为第五章误差建模的依据，理由是 \ph。
预报误差的分布特征（提前期—目标整点 MAE 热力图、配对折线、PIT 或 QQ 图、
残差自相关）汇总于图~\ref{fig:forecasterr}。

\begin{figure}[htbp]
  \centering
  \includefig[width=0.95\linewidth]{fig09_forecast_error.pdf}
  \caption{光伏预报误差结构（左上：提前期—目标整点 MAE 热力图；右上：同目标整点配对折线；
           左下：概率积分变换（PIT）或 QQ 图；右下：残差自相关函数。
           样本：附件 3 全部预报与附件 2 对应实际值配对，\(n=\ph\)）}
  \label{fig:forecasterr}
\end{figure}

\begin{guideblock}
图 \ref{fig:forecasterr} 的论证任务：\textbf{说明误差模型不是假设出来的，而是从数据里检验出来的}。
结论要写成"由数据可得：误差随提前期增大且存在系统性低估（bias \(=\ph\)），
故决策层需设置安全边际"，而不是只说"附件 3 精度较高"。
禁止话术："误差约 2.2 倍"（无配对设计与 \(n\)）——必须给配对比较、比值与样本量。
\end{guideblock}
